\documentclass[12pt, %
openright, 
oneside, %
%twoside, %TCC: Se seu texto tem mais de 100 páginas, descomente esta linha e comente a anterior
a4paper,    %
%english,   %
brazil]{facom-ufu-abntex2}

\usepackage{graphicx}
\graphicspath{{figuras/}{pictures/}{images/}{./}} % where to search for the images

\newcommand{\blue}[1]{\textcolor{blue}{#1}}
\newcommand{\red}[1]{\textcolor{red}{#1}}


\autor{Mateus Antônio Souza Silva} %TCC
\data{2026}
\orientador{Pedro Augusto Queiroz de Assis} %TCC
%\coorientador{Algum?} %TCC

% ---
% Informações de dados para CAPA e FOLHA DE ROSTO
% ---

\titulo{DESENVOLVIMENTO DE UM POSICIONADOR LINEAR} %TCC

\hypersetup{pdfkeywords={palavra 1}{palavra 2}{palavra 4}{palavra 4}{palavra 5}} %TCC

\begin{document} 
\frenchspacing 

% ----------------------------------------------------------
% ELEMENTOS PRÉ-TEXTUAIS
% ----------------------------------------------------------
%\pretextual
\imprimircapa
\imprimirfolhaderosto


% ---
% Inserir folha de aprovação
% ---
%
% \includepdf{folhadeaprovacao_final.pdf} %TCC: depois de aprovado o trabalho, descomente esta linha e comente o próximo bloco para incluir scan da folha de aprovação.
%
\begin{folhadeaprovacao}

  \begin{center}
    {\ABNTEXchapterfont\large\imprimirautor}

    \vspace*{\fill}\vspace*{\fill}
    {\ABNTEXchapterfont\bfseries\large\ CONSTRUÇÃO E CONTROLE DE UM POSICIONADOR LINEAR UTILIZANDO REALIMENTAÇÃO DE ESTADOS}
    \vspace*{\fill}
    
    \hspace{.45\textwidth}
    \begin{minipage}{.5\textwidth}
        \imprimirpreambulo
    \end{minipage}%
    \vspace*{\fill}
   \end{center}
    
   Trabalho aprovado. \imprimirlocal, 11 de março de 2026: %TCC:

   \assinatura{\textbf{\imprimirorientador} \\ Orientador}  
   \assinatura{\textbf{Jose Jean Paul Zanlucchi de Souza Tavares}}% \\ Convidado 1} %TCC:
   \assinatura{\textbf{Saint Clair Naves}}% \\ Convidado 2} %TCC:
   %\assinatura{\textbf{Professor} \\ Convidado 3}
   %\assinatura{\textbf{Professor} \\ Convidado 4}
      
   \begin{center}
    \vspace*{0.5cm}
    {\large\imprimirlocal}
    \par
    {\large\imprimirdata}
    \vspace*{1cm}
  \end{center}
  
\end{folhadeaprovacao}
% ---


%%As seções dedicatória, agradecimento e epígrafe não são obrigatórias.
%%Só as mantenha se achar pertinente.

% ---
% Dedicatória
% ---
%\begin{dedicatoria}
%   \vspace*{\fill}
%   \centering
%   \noindent
%   \textit{Dedico a \lipsum[10]}  %TCC:
%   \vspace*{\fill}
%\end{dedicatoria}
% ---

% ---
% Agradecimentos
% ---
%\begin{agradecimentos}
%Agradeço a \lipsum[30]. %TCC:
%\end{agradecimentos}
% ---

% ---
% Epígrafe
% ---
%\begin{epigrafe}
%    \vspace*{\fill}
%	\begin{flushright}
%		\textit{``Alguma citação que ache conveniente? \lipsum[10]''} %TCC:
%	\end{flushright}
%\end{epigrafe}
% ---

\begin{resumo} %TCC:
 

 \vspace{\onelineskip}
    
 \noindent
 \textbf{Palavras-chave}:  %TCC:
\end{resumo}



% ---
% inserir lista de ilustrações
% ---
\pdfbookmark[0]{\listfigurename}{lof}
\hypersetup{linkcolor = black}
\listoffigures*
\cleardoublepage
% ---

% ---
% inserir lista de tabelas
% ---
\pdfbookmark[0]{\listtablename}{lot}
\listoftables*
\cleardoublepage
% ---



% ---
\input{elementos/abreviaturas}

%% ---
%% inserir lista de símbolos, se for adequado ao trabalho. %TCC:
%% ---
%\begin{simbolos}
%  \item[$ \Gamma $] Letra grega Gama
%  \item[$ \Lambda $] Lambda
%  \item[$ \zeta $] Letra grega minúscula zeta
%  \item[$ \in $] Pertence
%\end{simbolos}
%% ---

% ---
% inserir o sumario
% ---
\pdfbookmark[0]{\contentsname}{toc}
\tableofcontents*
\cleardoublepage
% ---





% ----------------------------------------------------------
% ELEMENTOS TEXTUAIS
% ----------------------------------------------------------
\textual


% ----------------------------------------------------------
% Introdução
% ----------------------------------------------------------

\chapter[Introdução]{Introdução}
Existe uma demanda constante nas indústrias por melhoras na eficiência e na produção, seja via otimização de processos ou na automação dos mesmos, tendo cada vez mais a busca pelo controle automático de sistemas. Nesse cenário, sistemas de controle automático se tornam o alicerce para uma nova revolução de automação industrial \cite{bencomo2002control}. Em virtude disso, o domínio de técnicas de controle se torna uma necessidade na formação do engenheiro mecatrônico, visto a busca do mercado por profissionais engenheiros capacitados nas mais diversas técnicas de controle.

Sistemas de controle possuem duas linhas principais de pensamento \cite{kheir1996control}: uma é baseada na experiência prática enquanto a outra é baseada na teoria e matemática. Sendo assim, o engenheiro mecatrônico necessita compreender na prática os efeitos das soluções teóricas na resolução de problemas reais e a validação dos conceitos de modelagem, identificação de parâmetros e projeto de controladores. Assim, utilizar-se de protótipos didáticos em ambientes educacionais na área de controle é algo recomendável para uma formação consistente dos alunos \cite{silva2001midias}. Não se pode esquecer que a existência de protótipos didáticos pode ainda auxiliar como base para futuras pesquisas envolvendo controle digital, realimentação de estados ou identificação de parâmetros, seja pelos estudantes de graduação ou pelos de pós-graduação.

Neste contexto, este trabalho consiste no projeto e construção de uma bancada experimental de posicionamento linear acionada por motor CC, bem como o seu controle via realimentação de estados. O posicionador linear foi escolhido em virtude de possuir diversos parâmetros comuns e problemas tradicionais de controle \cite{netto2016relatorio}, sendo um modelo adequado para fins didáticos visto que não possui um nível elevado de complexidade. 

Assim, este trabalho tem como objetivo desenvolver e implementar um sistema de controle de posição para um posicionador linear, empregando um controlador por realimentação de estados embarcado em um microcontrolador. Busca-se construir uma bancada experimental didática, bem como modelar matematicamente seu sistema ao considerar os principais parâmetros físicos, criando a necessidade de identificar experimentalmente os parâmetros físicos do modelo. Será ainda feito o projeto de um controlador por realimentação de estados que atenda aos requisitos de desempenho definidos e sua implementação em hardware. Serão obtidos resultados experimentais com a bancada elaborada e resultados de simulação por meio do modelo matemático obtido, os resultados serão analisados e comparados, sendo discutindo as discrepâncias observadas entre os mesmos.

O restante do presente trabalho está estruturado do seguinte modo:

\begin{itemize}
    \item O Capítulo 2 descreve os componentes do posicionador linear em conjunto com sua arquitetura de funcionamento e modelagem matemática.
    \item O Capítulo 3 aborda os conceitos teóricos da realimentação de estados e desenvolve o controlador para o posicionador linear.
    \item No Capítulo 4 são apresentados os resultados obtidos e é realizada a análise desses. No capítulo são analisados os resultados de simulação e experimentais, sendo feita uma comparação entre os mesmos.
    \item O Capítulo 5 traz a conclusão do trabalho junto com sugestões de trabalhos futuros.
\end{itemize}

\include{capitulos/descricao_sistema}

\include{capitulos/descricao_controle}

\include{capitulos/resultados}

\chapter{Conclusões e sugestões de trabalhos futuros}

Em meio as constantes demandas do mercado por profissionais capacitados em técnicas de contole, tornou-se necessário introduzir a prática de controle aos estudantes de engenharia mecatrônica. Com essa linha de raciocínio, o objetivo deste trabalho era desenvolver uma bancada didática de um posicionador linear e implementar um sistema de controle de posição.

Para a realização desse objetivo, utilizou-se de um Arduino Mega para embarcar o controlador, tornando possível variar a tensão média aplicada ao motor CC via uma ponte H. Além disso, um modelo matemático no Espaço de Estados foi desenvolvido e validado, sendo realizados ensaios de resposta a um degrau de tensão e ensaios de torque do motor para obtenção dos coeficientes do modelo. Com a bancada construída e o modelo matemático validado, foi elaborado um controlador por Realimentação de Estados para realizar a ação de controle e escolhida uma referência de posição do tipo degrau, sendo colocado a teste via experimentos com a bancada e via simulação computacional.

Resultados de simulação demonstraram que o controlador era capaz de guiar o carro para a referência, embora não tenha antinjido os requisitos de projeto em virtude dos efeitos não lineares na simulação. Já os resultados experimentais demonstraram uma resposta do sistema mais lenta e mais próxima a um sistema de primeira ordem, mesmo assim, o controlador foi capaz de guiar o carro para a referência desejada e cumprir seu objetivo principal.

Por meio dos resultados, pôde-se constatar que as discrepâncias com os requisitos de projeto estipulados podem ser atribuídas a efeitos não lineares presentes no equipamento real, como a zona morta e o atrito seco, além é claro de incertezas dos parâmetros do modelo e comportamento não linear de algumas grandezas. 

Por fim, considerando as limitações do controlador escolhido e as dificuldades enfrentadas, possíveis trabalhos futuros podem se utilizar de técnicas mais avançadas de controle que levem em consideração as limitações físicas do equipamento real, como o Controle Preditivo Baseado em Modelo.

% ----------------------------------------------------------
% ELEMENTOS PÓS-TEXTUAIS
% ----------------------------------------------------------
\postextual


% ----------------------------------------------------------
% Referências bibliográficas
% ----------------------------------------------------------
\bibliography{bibliografia}


%% ----------------------------------------------------------
%% Apêndices TCC: só mantenha se for pertinente.
%% ----------------------------------------------------------

% ---
% Inicia os apêndices
% ---
\begin{apendicesenv}

% Imprime uma página indicando o início dos apêndices
\partapendices

\include{elementos/apendice_filtro}


\end{apendicesenv}
% ---

%---------------------------------------------------------------------
% INDICE REMISSIVO
%---------------------------------------------------------------------

\printindex



\end{document}